\documentclass[10pt,a4paper]{amsart}
\usepackage[margin=19mm]{geometry}
\usepackage{amsmath,amssymb,mathtools}
\usepackage[scr=rsfs,cal=euler]{mathalfa}
\usepackage{xfrac}
\usepackage[unicode,hidelinks,bookmarks=false]{hyperref}
\newcommand{\ie}{i.e.\ }
\newcommand{\cf}{cf.\ }
\newcommand{\resp}{respectively\ }
\newcommand{\Subsection}{\subsection}
\providecommand{\nobreakdash}{\nobreak\hbox{-}}
\setlength{\emergencystretch}{2em}
\multlinegap0pt
\usepackage{bm}
\usepackage{bbm}
\usepackage{dsfont}
\usepackage{xspace}
\usepackage{cancel}
\usepackage{mathtools}
\usepackage{amsthm}
\makeatletter
\@ifundefined{defi}{\newtheorem{defi}{Definition}}{}
\@ifundefined{lem}{\newtheorem{lem}[defi]{Lemma}}{}
\makeatother
\newcommand{\PP}{\mathrm{P}} %power set
\newcommand{\RR}{\mathbb{R}}
\newcommand{\TT}{\mathcal{T}} % thom functor
\newcommand{\eps}{\epsilon}
\title[Sinha's spectral sequence for long knots in codimension one ]{Sinha's spectral sequence for long knots in codimension one and non-formality of the little 2-disks operad}
\author{Syunji Moriya}
\date{Source: arXiv:2303.08111v3 (CC BY 4.0)}
\begin{document}
\maketitle

\noindent\emph{Source and licence.} Sinha's spectral sequence for long knots in codimension one and non-formality of the little 2-disks operad. Version arXiv:2303.08111v3 (CC BY 4.0), distributed under the Creative Commons Attribution 4.0 International licence, as recorded in the arXiv licence metadata for this exact version and shown on the abstract page https://arxiv.org/abs/2303.08111v3. Full rights notice: licenses/arxiv-2303.08111-CC-BY-4.0.txt.\par\smallskip

\noindent\textit{Editorial scope.} Counted unit: the Lem whose source label is \texttt{Ldiagonal\_bound} in the article (source0.tex lines 491--499, source label \texttt{Ldiagonal\_bound}), delivered together with the article's own complete original proof of it (source0.tex lines 500--551, one \texttt{proof} environment of 52 lines). No mathematics is written, repaired or re-derived: statement and proof are the article's own text line for line. Required context retained uncounted: Dpartition (lines 451--489). Every definition, display and label of those lines is retained unchanged. Omitted from this package and not redistributed: 1--450, 490--490, 552--2152. Nothing inside a retained excerpt is omitted or altered: every retained body is delivered exactly as the article's lines are written, so no omission receipt of any kind accompanies this package. Citations: the retained excerpts carry no citation command at all, so no bibliography entry of the article is transcribed and no reference is redistributed. Reference adaptation: none was needed and none was made; every reference inside the delivered unit resolves inside it, so no unresolved reference or citation is produced. Printed slips kept literal: no printed word, symbol, hypothesis or number of the counted statement or of its proof was altered. Layout and class: the article's own class is \texttt{amsart} with packages including setspace, amsthm, amscd, amsmath, amssymb, mathrsfs, picture, xy, here, enumerate; the wrapper is the lane's pinned \texttt{amsart} editorial wrapper at 10pt on A4, which restates the theorem-like environments the retained text needs and adds the article's own macro definitions that the retained text needs (\texttt{\textbackslash PP}, \texttt{\textbackslash RR}, \texttt{\textbackslash TT}, \texttt{\textbackslash eps}), with extra wrapper packages (bm, bbm, dsfont, xspace, cancel, mathtools). The delivered numbering is the unit's own. Assets: no retained span contains a figure environment, a graphics-inclusion command, a PDF-page inclusion, a table or a TikZ picture; no image file is copied into this package, the package declares no asset dependency and no image file is redistributed with it. Licence: version arXiv:2303.08111v3 of this article is distributed under the Creative Commons Attribution 4.0 International licence, as recorded in the arXiv licence metadata for this exact version and shown on the abstract page https://arxiv.org/abs/2303.08111v3. A full rights notice is delivered at licenses/arxiv-2303.08111-CC-BY-4.0.txt. Characters: every retained excerpt is pure ASCII, so the delivered engine, XeLaTeX with its default Latin Modern text font, reports no missing character.
\begin{defi}\label{Dpartition} Let $P\in \PP_n$ be a partition and set $p=\# P-2$.
\begin{enumerate}




\item Let $\nu_{P}$ be the $\eps_P$-neighborhood of $e_P(\RR^{dp})$ i.e.
\[
\nu_P=\{y\in \RR^{dn}\mid |y-e_P(x)|<\eps_P \text{ for some } x\in \RR^{dp}\}.
\]  For  $\alpha, \beta\in P^\circ$,  we define a   subspace $D_{\alpha\beta}$ of $\RR^{dp}=(\RR^d)^p$ by
\[
D_{\alpha\beta}=D_{\alpha\beta}(P)=\{(x_\gamma)_{\gamma \in P^\circ} \mid |x_\alpha-x_\beta|\leq d_{\alpha\beta}(P)\}, 
\]
where 
\[
d_{\alpha\beta}(P)=\rho c_{\alpha \beta}-\eps_P.
\]
We denote by $E_{\alpha}=E_\alpha(P)$  the subspace of  elements $(x_\gamma)_\gamma$ satisfying the following condition:
\[
\begin{split}
|x_\alpha|& \geq 1-\rho c_\alpha/2+\eps_P,\  \text{or} \\
  x_\alpha&\leq_1\, -1+\rho c_{\leq \alpha}-\eps_P,\  \text{or} \\  x_\alpha&\geq_1\,  1-\rho c_{\geq \alpha}+\eps_P.
\end{split}
\] 
Set 
\[
E_P=\bigcup_{\alpha\in P^\circ}E_\alpha.
%\qquad D_P=E_P\cup\bigcup_{\alpha,\beta\in P^\circ, \alpha<\beta}D_{\alpha \beta}.
\] 
Let $\pi_P:\RR^{dn}\to e_P(\RR^{dp})=\RR^{dp}$ be the orthogonal projection. 
%We regard $\nu_P$ as a disk bundle with the restriction of $\pi_P$ and set $\nu_P|_X= \nu_P\cap \pi_P^{-1}(X)$.
We set
\[
\TT_{\emptyset_P}=\RR^{dn}/\{(\RR^{dn}-\nu_P)\cup \pi_P^{-1}(E_P)\}.
\]
 For two pieces $\alpha, \beta\in P^\circ$ with $\alpha<\beta$, let $\TT_{\alpha\beta}$ denote the subspace of $\TT_{\emptyset_P}$ consisting of the basepoint and the elements represented by  $x\in \nu_P$ satisfying $\pi_P(x)\in D_{\alpha\beta}$.

\end{enumerate}
\end{defi}

\begin{lem}\label{Ldiagonal_bound}
Let $Q$ be a subdivision of a partition $P\in \PP_n$. We have
\[
(\RR^{dn}-\nu_Q)\cup \pi_Q^{-1}(E_Q)\subset (\RR^{dn}-\nu_P)\cup \pi_P^{-1}(E_P).
\]
In particular, the identity on $\RR^{dn}$ induces the collapsing map 
$\delta'_{P,Q}:\TT_{\emptyset_Q}\to \TT_{\emptyset_P}.
$
\end{lem}

\begin{proof}
 Let $y\in \RR^{dn}$.
Write
\[
\begin{split}
\pi_{Q}(y)&=(x_\gamma)_{\gamma\in Q^\circ}, \\
\pi_{P}(y)&=(\bar x_{\gamma'})_{\gamma'\in P^\circ},\\
e_{P,Q}(\pi_P(y))&=(x'_\gamma)_{\gamma\in Q^\circ} .
\end{split}
\] Suppose $y\in \nu_P$.  Since the image of $e_P$ is contained in the image of $e_Q$ and the map $\pi_{Q}$ sends $y$ to its closest point in the image of $e_Q$, we have
\begin{equation}
|y-e_Q(x_\gamma)|\leq |y-e_Q(x'_\gamma)|=|y-e_P(\pi_P(y))|< \eps_P\ (<\eps_Q). \label{Eemb}
\end{equation}
So we see $y\in \nu_Q$. This means $\RR^{dn}-\nu_Q\subset \RR^{dn}-\nu_P$. We shall show $\nu_Q\cap \pi_Q^{-1}(E_Q)\subset (\RR^{dn}-\nu_P)\cup \pi_P^{-1}(E_P)$. Let $\alpha\in Q^\circ$ be a piece and $\alpha'$  the piece of $P$ including $\alpha$. If $y\not\in \nu_P$, we have nothing to prove, so suppose  $y\in \nu_P\cap\pi_Q^{-1}(E_Q)$ and we shall prove $y\in \pi_P^{-1}(E_P)$. Intuitively speaking, $x_\alpha$ and $\bar x_{\alpha'}$ are different in general but the difference is sufficiently small since we have chosen the radius of the tubular neighborhood $\nu_P$ sufficiently small. Since $\eps_P<\eps_Q$ and $\alpha\subset \alpha'$, the range given by the inequalities defining $E_{\alpha'}(P)$ is wider than the one given by the inequalities of $E_\alpha(Q)$ and the margin covers the difference between $x_\alpha$ and $\bar x_{\alpha'}$ so we have $y\in \pi_P^{-1}(E_P)$. We shall give a rigorous proof.  Suppose further $|x_\alpha|\geq 1-\rho c_\alpha/2+\eps_Q$. 
By the inequality (\ref{Eemb}), we have
\[
|e_Q(x'_\gamma)-e_Q(x_\gamma)|\leq |e_Q(x'_\gamma)-y|+|y-e_Q(x_\gamma)|\leq 2\eps_P.
\]
As $e_Q$ is a composition of a diagonal map with a parallel transport, we have
\begin{equation}
|(x'_\gamma)_\gamma-(x_\gamma)_\gamma |\leq 2\eps_P \quad (\text{so} \ |x'_\gamma-x_\gamma|\leq 2\eps_P\ \text{for each}\ \gamma\in Q^\circ). \label{Edifference}
\end{equation}
Let $\beta\subset \alpha'$ be the  set of elements smaller than the minimum of $\alpha$. We easily see 
\begin{equation}
\bar x_{\alpha'}-x'_\alpha=\frac{\rho}{2}(c_{\alpha'}-2c_\beta-c_\alpha)u. \label{Earrange}
\end{equation}
Putting these (in)equalities together, we see
\[
\begin{split}
|\bar x_{\alpha'}|& \geq |x_\alpha|-|x_\alpha-x'_\alpha|-|x'_\alpha-\bar x_{\alpha'}| \\
& \geq 1-\rho c_\alpha/2+\eps_Q-2\eps_P-\rho(c_{\alpha'}-c_\alpha)/2 \\
& \geq 1-\rho c_{\alpha'}/2+\eps_P
\end{split}
\]
since we have 
\[
\eps_Q-2\eps_P=\frac{1-2\cdot 8^{p-q}}{8^{n-q}}\geq \frac{1-2/8}{8^{n-q}}>\eps_P.
\]
We have shown that the first inequality in Definition \ref{Dpartition}(2) for $Q$ and the condition $y\in \nu_P$ imply the corresponding inequality for $P$. 
Similarly, suppose $x_\alpha\leq_1 -1+\rho c_{\leq \alpha}-\eps_Q$ (and $y\in \nu_P$). Let $\alpha'$ be  as above.
We see 
\[
\begin{split}
\bar x_{\alpha'} & =\bar x_{\alpha'}-x'_\alpha+x'_\alpha-x_\alpha+x_\alpha \\
& \leq_1 \frac{\rho}{2}(c_{\alpha'}-2c_\beta-c_\alpha)+2\eps_P+(-1+\rho c_{\leq \alpha}-\eps_Q) \\
&=-1+\rho c_{\leq \alpha'}-(\eps_Q-2\eps_P)  \quad (\because \ c_{\leq \alpha}+(c_{\alpha'}-2c_\beta-c_\alpha)/2=c_{\leq \alpha'}) \\
&<-1+\rho c_{\leq \alpha'}-\eps_P.
\end{split}
\]
This is  the second inequality in Definition \ref{Dpartition}(2) for $P$. 
Similarly, we see $\bar x_{\alpha'}\geq_1 1-\rho c_{\geq \alpha'}+\eps_P$ if $ x_{\alpha}\geq_1 1-\rho c_{\geq \alpha}+\eps_P$ and $y\in \nu_P$. Thus, we have shown the claimed inclusion.
\end{proof}

\end{document}
